% FDCA v1.0.1 -- unreleased GitHub accuracy revision candidate.
% Original v1.0.0 source and archival PDF are preserved separately.
% CPU-only reanalysis; no new image generation, model training or GPU inference.
\documentclass[11pt,a4paper]{article}
\PassOptionsToPackage{colorlinks=true,linkcolor=black,citecolor=black,urlcolor=blue}{hyperref}
\usepackage[T1]{fontenc}
\usepackage[utf8]{inputenc}
\usepackage{amsmath,amssymb,amsthm}
\usepackage{newtxtext,newtxmath}
\usepackage[scaled=0.9]{helvet}
\usepackage[margin=1in]{geometry}
\usepackage{booktabs,array,longtable}
\usepackage{graphicx,caption,subcaption,placeins}
\usepackage{enumitem,xcolor,microtype}
\usepackage{natbib,fancyhdr,hyperref,xurl}
\setlength{\emergencystretch}{2em}
\hypersetup{pdftitle={First-Divergence Consequence Analysis: Separating Quantization Onset and Output Discrepancy in Masked Visual Generators},pdfauthor={Munsik Kim},pdfsubject={Unreleased scientific accuracy revision},pdfkeywords={masked visual generators, quantization, coupling, paired block inference},pdfcreator={TeX; FDCA v1.0.1 revision candidate}}
\bibliographystyle{plainnat}
\theoremstyle{plain}
\newtheorem{theorem}{Theorem}[section]
\newtheorem{proposition}[theorem]{Proposition}
\theoremstyle{definition}
\newtheorem{definition}[theorem]{Definition}
\theoremstyle{remark}
\newtheorem{remark}[theorem]{Remark}
\newcommand{\TV}{\mathrm{TV}}
\newcommand{\Ind}[1]{\mathbf{1}\{#1\}}
\newcommand{\Ex}{\mathbb{E}}
\newcommand{\Prob}{\mathbb{P}}
\newcommand{\Mask}{\texttt{M}}
\newcommand{\diam}{\mathrm{diam}}
\newcommand{\rholabel}{\widetilde{\rho}}
\newcommand{\rhoatt}{\rho_{\mathrm{att}}}
\newcommand{\Meq}{M_{\mathrm{equal}}}
\newcommand{\Msplit}{M_{\mathrm{split}}}
\newcommand{\revisiongenerated}{../../generated/v1.0.1}
\setlength{\headheight}{15pt}
\pagestyle{fancy}
\fancyhf{}
\fancyhead[L]{\small First-Divergence Consequence Analysis}
\fancyhead[R]{\small v1.0.1 revision candidate}
\fancyfoot[C]{\thepage}
\title{\bfseries First-Divergence Consequence Analysis:\\
Separating Quantization Onset and Output Discrepancy\\
in Masked Visual Generators}
\author{Munsik Kim\\\small Independent Researcher, Seoul, Republic of Korea\\
\small ORCID: \href{https://orcid.org/0009-0008-9350-2435}{0009-0008-9350-2435}}
\date{Version 1.0.1 --- unreleased revision candidate\\[2pt]
\small Accuracy revision prepared 29 October 2026. Technical preprint; not peer reviewed.}
\begin{document}
\maketitle
\thispagestyle{fancy}

\begin{abstract}
\noindent
A small local change in an iterative generator can cause either no sampled-token
change or a mismatch that propagates to the output. First-Divergence Consequence
Analysis (FDCA) separates these events through an explicit measurement contract:
a common anchor state, a natural quantization seed drawn under a declared
coordinatewise coupling, a common propagation kernel, audited random events,
and inference over class--generation-seed blocks. The mathematical contribution
consists of a stable-set additive-cost bound, an expectation identity for
first-divergence paths, and a finite-horizon influence envelope under global
hybrid-state assumptions. These results do not prescribe opposite trends in
incidence and consequence and do not automatically bound LPIPS.

Under a fixed 32-step Halton schedule we compare one FP32-versus-W8 transition
and propagate both branches with a common W8 suffix. In separate H-MaskGIT-T
($n=48$) and H-MaskGIT-S ($n=40$) CFG3 panels, moving from target unresolved
fraction $0.3$ to $0.5$ lowers batch split probability while increasing both the
equal-block split-conditional LPIPS and the mean of blockwise incidence-weighted
LPIPS. Original paired-block intervals support conditional-LPIPS increases of
$0.02000$ and $0.02038$, and incidence-weighted increases of $0.001714$ and
$0.001966$. The latter is a mean of products, not a product of panel means.
A labeled post-hoc CPU reanalysis checks aggregation, bootstrap robustness and
schedule dependence. The commit batch shrinks from 14 to 8 positions; mean
coordinate TV does not show the same decrease as batch split probability.

The contribution is an auditable separation of onset and conditional output
discrepancy in this intervention. The observations establish neither a universal
reverse-time law, persistent-quantization quality degradation, human preference,
nor a validated precision-selection policy. Frozen source resolves per-output-channel quantization, including a documented
deviation from the per-tensor plan; public compact evidence retains further
reproducibility limits.
\end{abstract}

\noindent\fbox{\parbox{0.96\textwidth}{\small
\textbf{Revision and evidence boundary.} This candidate corrects the interpretation
of existing measurements and adds CPU reanalysis of public compact tables.
The v1.0.0 PDF, source, preregistrations and frozen results remain historical
records. No GPU inference, model training, new images or new perceptual measurements
were performed. The DOI \href{https://doi.org/10.5281/zenodo.21965747}{10.5281/zenodo.21965747}
identifies archived v1.0.0, not this unreleased revision.}}

\section{Introduction}
Local approximation diagnostics and endpoint metrics answer different questions.
Logit total variation (TV) measures a change in a categorical law. A paired image
distance measures discrepancy between two specified outputs. A coupling connects
them by declaring a joint sampling experiment, including when the outputs remain
equal and what happens after the first sampled-token mismatch.

FDCA makes that joint experiment explicit for fixed monotone masked decoding.
At a common anchor, the reference and approximate categorical laws produce a
natural seed under coordinatewise maximal coupling. Both resulting states then
use the same approximate kernel. This separates the probability of a batch split
from the output consequence conditional on that split. Event-addressed random
numbers, zero controls and block-level aggregation make the contract inspectable.
The construction addresses a measurement question; it does not establish the
optimal coupling for an image metric or the best deployment policy.

First disagreement is a classical stopping-time construction
\citep{lindvall2002coupling,vollering2016maximal,ernst2019mexit}.

Single-point sensitivity interventions are also not unique to FDCA. In particular,
Q-ARVD~\citep{tang2026qarvd} studies frame-specific quantization sensitivity
followed by a full-precision suffix; FDCA's natural-seed panels instead use a common W8 suffix. The distinction
lies in the declared discrete onset law, conditional consequence, natural seed,
coupling and RNG audits, and paired block inference together, rather than in the
general observation that an early error can propagate.

The strongest empirical result is bounded but reproducible. Between target
unresolved fractions $0.3$ and $0.5$, two CFG3 W8 checkpoint panels show
$H\downarrow$, $\Meq\uparrow$, and $R\uparrow$, where $H$ is mean batch split
probability, $\Meq$ is the equal-block conditional LPIPS, and $R$ is the mean
blockwise product. Batch size, state, coordinates, remaining horizon and seed
composition all vary across anchors. The result therefore motivates measuring
incidence and consequence separately; it does not identify a unique cause or
imply the pattern holds at every anchor.

\paragraph{Contributions.}
\begin{enumerate}[leftmargin=1.6em,itemsep=3pt]
\item \textbf{A specified measurement framework.} A common anchor, one natural
FP32-versus-W8 seed transition, common W8 propagation, coordinatewise coupling,
semantic RNG and class--seed block inference define the measured discrepancy.
\item \textbf{Finite monotone-decoder theory.} We give a pathwise stable-set
bound, a first-divergence expectation identity and a clipped DAG envelope.
The envelope requires global state classes containing both branches and the
hybrids used in the proof. Its finite inverse uses nilpotency, not contraction.
\item \textbf{Two checkpoint panels with a supported contrast.} In the stated
CFG3 W8 intervention, lower batch split probability at target $0.5$ coexists
with higher equal-block conditional and incidence-weighted LPIPS than at $0.3$.
Matched timing, CFG1 and saturated W4 panels have distinct supporting roles.
\item \textbf{An accuracy and reproducibility correction.} We distinguish
$\Ex_b[hm]$, $\Ex_b[h]\Ex_b[m]$, and the split-weighted population conditional
mean; reproduce legacy numerical checks; and provide labeled post-hoc diagnostics,
input hashes, generated tables and figures, and a buildable revised manuscript.
\end{enumerate}

\section{Related Work}
\label{sec:related}
\subsection{Masked generation and schedule dependence}
MaskGIT introduced iterative bidirectional visual-token generation
\citep{chang2022maskgit}. Halton scheduling uses a spatially dispersed sequence
\citep{halton1960efficiency,besnier2025halton}.
Work on factorized masked-diffusion approximations and inference schedules
analyzes bias caused by simultaneous unmasking
\citep{park2025jump,lavenant2025error,chen2026optimal}.
The MaskGIT sampler analysis of \citet{hayakawa2025demystifying} studies its
implicit temperature and order-selection mechanism. These results make the
schedule part of the model of the experiment, not a neutral time label.
FDCA fixes a branch-independent schedule and studies the declared paired
implementation intervention; changing the schedule requires a new contract.

\subsection{Coupling and comparison arguments}
Maximal agreement and MEXIT ask how long stochastic processes can remain equal
\citep{vollering2016maximal,ernst2019mexit}.
Coordinatewise maximal coupling minimizes one coordinate's mismatch probability.
It does not generally maximize whole-batch agreement or minimize a terminal
LPIPS distance. Dobrushin comparison arguments bound discrepancy through local
conditional perturbations \citep{rebeschini2014comparison}. Our finite,
time-ordered influence matrix follows this lineage. The monotone commit order
makes it nilpotent; measured neural local profiles remain different from the
global suprema required for a certificate.

\subsection{Visual quantization and error propagation}
Q-ARVD \citep{tang2026qarvd} quantizes one autoregressive video frame
while keeping the prefix and subsequent frames full precision in its sensitivity
experiment (Section 3.1, Equation 3). This is a close precedent for localized
shock/propagation analysis. FDCA instead defines a categorical natural-seed
split and propagates both branches using a common W8 suffix. Neither a
one-time perturbation nor the observation of error propagation is an exclusive
novelty claim here.

For continuous diffusion, \citet{liu2026propagation} analyze accumulated
quantization error through a first-order propagation decomposition and develop
TCEC compensation. That local linearization and continuous denoising setting
differ from FDCA's categorical split law and its monotone discrete-state
contract. PTQ4ARVG \citep{liu2026ptq4arvg} develops quantization/calibration
methods for autoregressive visual generation, and Shift-and-Sum
\citep{moon2026shift} targets attention-value reconstruction and calibration
across scales. Caching and mixed quantization also seek computational savings
\citep{yan2025lazymar,wu2025quantcache}.

FDCA's contribution is the combined onset/consequence measurement contract,
not demonstrated superiority over these quantizers or sensitivity methods.
The present data do not show that it predicts outcomes or selects precision
better than logit TV, seed geometry, or simple early protection. Those comparisons
require held-out policy tests (Section~\ref{sec:future}).

\subsection{Output metrics}
LPIPS \citep{zhang2018unreasonable} and DINOv2 embedding distance
\citep{oquab2024dinov2} are paired output discrepancies. A ResNet-50 classifier
\citep{he2016deep} supplies secondary diagnostics. A greater distance between
the two specified outputs establishes neither a consistent quality loss from a
reference distribution nor human detection or preference.

\section{Setup, Estimands and Aggregation}
\label{sec:setup}
\subsection{Fixed monotone masked decoder}
Let $I$ be a finite set of $N$ token positions, $V$ a finite codebook, and
$\Mask\notin V$. A branch-independent schedule is an ordered partition
$\mathcal C=(C_1,\ldots,C_K)$ of $I$, fixed independently of either branch's
realized tokens. Positions in $C_{k+1}$ are committed at transition $k+1$.
Committed token values are never revised. At a pre-transition state $x_k$,
reference and approximate batch laws factorize as
\[
 P_k(\cdot\mid x_k)=\bigotimes_{u\in C_{k+1}}p_{k,u}(\cdot\mid x_k),\qquad
 Q_k(\cdot\mid x_k)=\bigotimes_{u\in C_{k+1}}q_{k,u}(\cdot\mid x_k).
\]
This product assumption, with independent coordinate coupling events conditional
on the state, is needed for the split-probability product formula. It is not
needed for the stable-set argument. Any additional memory affecting transitions
must be part of the state; token grids alone are insufficient for a model with
branch-dependent hidden memory or caches.

\begin{definition}[First divergence]
For two coupled processes starting equal, let
$\tau=\inf\{k\geq1:X_k\neq Y_k\}$, with $\tau=\dagger$ if no split occurs.
We use $G_k=\{\tau>k\}$ for survival without a prior split, treating
$\tau=\dagger$ as surviving every finite stage.
\end{definition}
The terminal Hamming cost is
$d_\lambda(x,y)=\sum_u\lambda_u\Ind{x_u\neq y_u}$, where
$\lambda_u\geq0$ and $\sum_u\lambda_u=1$.

\subsection{Unresolved fraction and schedule}
\label{sec:rho}
The label $\rholabel$ is a \emph{target unresolved-token fraction}, not a
quantization strength or a generation-step number. Larger fractions occur
earlier in decoding. The nearest audited pre-transition state has attained
fraction $\rhoatt$, which is generally different from the target.
\begin{table}[tb]
\centering\small
\caption{Frozen 32-step, $16\times16$ Halton schedule. Transition numbers are
one-indexed. The target $0.1$ anchor is the final flush of 41 remaining tokens.}
\label{tab:anchors}
\input{\revisiongenerated/schedule.tex}
\end{table}

\subsection{Block estimands and population conditional means}
\label{sec:aggregation}
A block $b$ denotes a class--generation-seed pair. Within one panel and anchor
$r$, let $S$ denote a split of the single anchor transition and define
\begin{align}
 h_{b,r}&=1-\prod_{u\in C(r)}[1-\TV(p_u,q_u)],\label{eq:incidence}\\
 m_{b,r}&=\Ex_\Gamma[d\mid S,b,r], & R_{b,r}&=h_{b,r}m_{b,r}.
\end{align}
The value $h_{b,r}$ is the conditional batch split probability under the specified
coordinate-independent coupling at the common state. It is computed from the
categorical laws (up to floating-point evaluation), not from a count of sampled
split events. Its across-block
uncertainty is therefore not a binomial Wilson interval. The conditional mean
$m_{b,r}$ is estimated from repeated conditional rollouts and retains Monte Carlo
error. For LPIPS, $d$ denotes the distance between the two decoded terminal grids.

For the population that first chooses a block with equal weight, define
\begin{equation}
 H=\Ex_b h_{b,r},\qquad \Meq=\Ex_b m_{b,r},\qquad
 R=\Ex_b[h_{b,r}m_{b,r}],\qquad \Msplit=R/H\quad(H>0).
 \label{eq:aggregates}
\end{equation}
Here $\Meq$ is the \emph{equal-block} conditional mean. In contrast, $\Msplit$
is the conditional mean after drawing an equal-weight block and conditioning
that population on a split; this reweights blocks in proportion to $h$.
The historical tables called $\Meq$ ``conditional LPIPS.'' Throughout this
revision that phrase is qualified as equal-block conditional LPIPS.
In general,
\begin{equation}
 R=H\Meq+\operatorname{Cov}_b(h,m),\qquad
 \operatorname{Cov}_b(h,m)=n^{-1}\sum_b(h_b-H)(m_b-\Meq).
 \label{eq:covariance}
\end{equation}
Thus multiplying printed panel means is not the published incidence-weighted
estimator. The covariance uses population normalization $1/n$.

If $h_{b,r}=0$, its contribution to $R$ is zero, but a conditional mean on that
zero-probability event is unobserved, not an observed zero. An all-block
$\Meq$ then requires an explicit domain convention and cannot be manufactured
by zero filling; report the positive-$h$ domain instead. If $H=0$, $\Msplit$
is undefined. All reported primary data have positive $h$.

\subsection{One shock and persistent approximation are different estimands}
\label{sec:estimands}
The natural-seed panels use a common FP32 prefix state, one FP32-versus-W8
transition, and a common W8 kernel on both subsequent branches. Equal branches
remain equal under the shared self-coupling. Conditional on a seed split,
the suffix measures state-mediated propagation without fresh same-state
FP32-versus-W8 innovation. Persistent approximation would instead compare an
FP32 branch and a W8 branch throughout the suffix and include further innovation.
The current $h_{b,r}$ values do not condition on survival along those persistent
paths. Summing them or multiplying their complements across anchors does not
estimate a persistent first-divergence CDF or its total expected discrepancy.

\section{Theory and Its Scope}
\label{sec:theory}
\subsection{Stable-set additive-cost bound}
\begin{definition}[Future-stable set]
At a common state $s$ before transition $k+1$, $F_k(s)$ contains coordinates equal
in both branches that no allowed continuation of either branch revises.
\end{definition}
\begin{theorem}[Stable-set bound]
\label{thm:stableset}
Assume the finite common monotone schedule and equal initial states above.
Let $d(x,y)=\sum_u\lambda_uc_u(x_u,y_u)$ with finite nonnegative
$\lambda_u$, nonnegative coordinate costs, $c_u(a,a)=0$, and bounded
$\diam(c_u)=\sup_{a,b}c_u(a,b)<\infty$ over the terminal domain.
Then on $\{\tau\neq\dagger\}$,
\[
 d(X_K,Y_K)\leq \sum_{u\notin F_{\tau-1}(X_{\tau-1})}\lambda_u\diam(c_u),
\]
and the terminal cost is zero on $\{\tau=\dagger\}$. Taking expectations
gives the corresponding expected-cost bound.
\end{theorem}
\begin{proof}
Immediately before the first split both states agree. Every coordinate in the
future-stable set remains equal through all allowed continuations, so its
terminal cost is zero. Bound each remaining nonnegative coordinate cost by its
finite diameter, multiply by its nonnegative weight, and sum. On no split,
all terminal coordinates agree and every diagonal cost is zero.
\end{proof}
For uniform Hamming cost this is the fraction of positions still capable of
changing. It is often loose early in decoding. Signed costs and negative weights
are outside the stated theorem. LPIPS is not a coordinate-additive token cost;
a decoder-to-metric bound with valid constants would be needed to transfer this
bound to LPIPS, and none is established here.

\subsection{Full-path expectation identity}
\label{sec:decomp}
For a specified coupled \emph{full-path} law, define
\[
 h_k(s)=\Prob(\tau=k+1\mid G_k,X_k=Y_k=s),\quad
 m_k(s)=\Ex[d(X_K,Y_K)\mid G_k,X_k=Y_k=s,\tau=k+1].
\]
\begin{proposition}[Onset--consequence identity]
\label{prop:decomp}
For an integrable terminal cost with zero cost on $\{\tau=\dagger\}$,
\[
 \Ex[d(X_K,Y_K)]=\sum_{k=0}^{K-1}
 \Ex[\Ind{G_k}h_k(X_k)m_k(X_k)].
\]
Zero-hazard terms contribute zero without assigning an observed conditional mean.
\end{proposition}
\begin{proof}
The disjoint events $\{\tau=k+1\}$ partition the event of a split. Expand the
expected cost over them and apply conditional expectation first on $G_k$ and the
common state, then on the split event. The no-split contribution is zero by
assumption.
\end{proof}
This is the law of total expectation. It implies no sign relationship between
$h$ and $m$. The one-anchor single-shock identity $R_{b,r}=h_{b,r}m_{b,r}$
is a different specialization: it has one allowed innovation transition,
not a sequence of persistent hazards. Mixing the anchor interventions does not
recover the survival-conditioned quantities of the proposition.

\subsection{Clipped innovation--influence DAG envelope}
At stage $k$, let $\mathcal S_k$ contain every possible state of both branches
and every intermediate state obtained by changing previously committed
coordinates one at a time between them. Unresolved coordinates and the stage
are fixed consistently in this class. For $u\in C_{k+1}$, define global suprema
\[
 s_u=\sup_{x\in\mathcal S_k}\TV(p_{k,u}(\cdot\mid x),q_{k,u}(\cdot\mid x)),
\qquad
 M_{uv}=\sup_{\substack{x,y\in\mathcal S_k\\x\sim_v y}}
 \TV(q_{k,u}(\cdot\mid x),q_{k,u}(\cdot\mid y)),
\]
where $v$ is in a strictly earlier commit block and $x\sim_v y$ means equality
except at $v$. Set $M_{uv}=0$ for other pairs.
\begin{theorem}[Clipped DAG envelope]
\label{thm:dag}
Assume the fixed branch-independent monotone schedule, batch laws determined
by the token grid and stage, compatible product marginals, global hybrid-state
coverage above, and conditional coordinatewise maximal
coupling of $p_{k,u}(\cdot\mid X_k)$ with $q_{k,u}(\cdot\mid Y_k)$ at every
commit. In commit order set
\[
 \bar r_u=\min\{1,s_u+\sum_{v\prec u}M_{uv}\bar r_v\}.
\]
Then $\Prob(X_K[u]\neq Y_K[u])\leq\bar r_u$ and
$\Ex[d_\lambda(X_K,Y_K)]\leq\lambda^\top\bar r$.
\end{theorem}
\begin{proof}
Write $\pi_u$ for terminal mismatch probability at $u$. Monotonicity makes it
equal to mismatch probability at that coordinate's commit. Maximal coupling
and the TV triangle inequality give a same-state term at most $s_u$ and a
state-change term. Interpolate from $X_k$ to $Y_k$ through the assumed hybrids.
Changing an earlier mismatched coordinate $v$ contributes at most $M_{uv}$;
a coordinate that agrees contributes zero. Taking expectations bounds the
state term by $\sum_{v\prec u}M_{uv}\pi_v$. Since probabilities are at most
one, $\pi_u\leq\min\{1,s_u+\sum_{v\prec u}M_{uv}\pi_v\}$.
Induction in commit order gives the asserted coordinate bounds. Sum with
nonnegative Hamming weights to obtain the terminal bound.
\end{proof}
The matrix is strictly block-triangular with $K$ commit blocks, so $M^K=0$ and
$(I-M)^{-1}=\sum_{j=0}^{K-1}M^j$. No contraction or spectral-radius assumption
is required. This is a finite-horizon statement, not infinite-time stability.

\begin{remark}[Measured profiles and metric transfer]
\label{rem:notcert}
Finite local measurements, q95 profiles and calibration maxima are not global
supremum certificates. T0 enumerates the finite synthetic state classes; it
does not validate these assumptions for arbitrary neural states. P05R's neural
envelope was empirically vacuous and P1M found limited incremental within-anchor
prediction. Neither these profiles nor observed LPIPS trends prove a neural
LPIPS bound. Confidence samplers that rewrite token values, continuous diffusion,
and new architectures require separate assumptions and definitions.
\end{remark}

\section{Implementation and Measurement Contract}
\label{sec:contract}
\subsection{Generator and state}
The official Halton-MaskGIT source is pinned to
\nolinkurl{f61b0a1314717004dc7487531fd16a8bb71e1888}
\citep{besnier2025halton}. H-MaskGIT-T uses hidden width 384, depth 6 and six
attention heads; H-MaskGIT-S uses width 512, depth 12 and six heads.
Both have a $16\times16$ grid and use 32 Halton steps with
\texttt{randomize=False}. The Small asset ledger pins
\nolinkurl{ffc7fd4c5fc7b16010acc4aa342310484d7de62a},
832,274,338 bytes and SHA-256
\nolinkurl{760b9937d16157e7e518fbf405efe281552933242b552d5e8262cb10507fb0e2}.
These are recorded provenance facts; the compact CPU analysis does not load
or independently rehash that external checkpoint.
The confidence sampler was inspected and excluded from the theorem because
committed values can be rewritten even when mask occupancy is monotone.

\subsection{Approximation provenance}
\label{sec:quantizer}
Frozen runner snapshots identify deterministic symmetric
\emph{per-output-channel} round-to-nearest (RTN) quantization of eligible
transformer linear weights. The imported \texttt{LinearWeightHomotopy} class
constructs W4/W8 weights with \texttt{rtn\_per\_output\_channel}; the function
reduces the weight maximum over input dimension 1 while retaining that dimension,
so its scale tensor has shape $[\mathrm{out},1]$:
\[
 Q_b(W_{o,:})=\operatorname{clip}\!\left(
 \operatorname{round}(W_{o,:}/\Delta_o),-q_{\max},q_{\max}\right)\Delta_o,
 \qquad \Delta_o=\max_j|W_{o,j}|/q_{\max}.
\]
All-zero rows return zero directly. PyTorch deterministic \texttt{round}
uses ties to even and the zero point is zero. W8 has $q_{\max}=127$.
These are weight-only fake-quantized weights used in FP32 computation;
activations, norms, embeddings, biases and the VQ decoder remain FP32.

The historical P2P preregistration instead names
\texttt{per\_tensor\_symmetric\_RTN}. This is a preserved plan-to-source
deviation; its YAML is not silently rewritten. Matching frozen P2P/P3S runner
imports and code-manifest hashes to the public quantizer source resolves the
source-level granularity. This revision does not load models or independently
regenerate historical runtime scale tensors or quantized-weight hashes. The
sanitized receipt is \nolinkurl{docs/evidence/QUANTIZER_SOURCE_RECEIPT.json}.
No packed-kernel runtime, speedup or VRAM saving is measured by this
output-fidelity experiment. W4 is retained only as saturated stress evidence.

\subsection{Coordinate coupling and split law}
For categorical laws $p,q$, draw $x\sim p$, accept $y=x$ with probability
$\min\{1,q(x)/p(x)\}$, and otherwise draw from $(q-p)_+/\TV(p,q)$.
If $\TV(p,q)=0$, return $(x,x)$ without normalizing a zero residual.
Then $X\sim p$, $Y\sim q$, and $\Prob(X\neq Y)=\TV(p,q)$.
Independently addressed coordinate events yield Equation~\eqref{eq:incidence}
for a factorized commit batch conditional on its common state. Coordinate
maximality is not batch, path or LPIPS optimality. The coupling choice changes
the paired estimand, even when both marginal laws are unchanged.

\subsection{RNG and logging}
Semantic keys identify experiment, cell, block, anchor, replicate, attempt,
step, position, event kind and draw. Separate event kinds address reference
proposals, accept/reject events, residual draws and suffix events. Registries
classify intentional shared reference streams and reject unclassified duplicate
or missing keys. P2B retains equal as well as mismatched selected-token pairs;
the historical reports find no selected MASK or clamp-induced mismatch.
Zero-seed controls check equal terminal grids, images and metrics. These are
historical execution/audit findings, not newly rerun generation in this revision.

LPIPS uses the AlexNet backbone, DINOv2 uses ViT-S/14 cosine distance, and
secondary P2P diagnostics include classifier JSD, SSIM and pixel $L_1$.
A smoke numerical audit selected scalar metric inference after a scalar-versus-
batch tolerance failure. The compact data retain scores; they do not independently
recreate checkpoint loading, image decoding, feature extraction or human ratings.

\section{Experimental Design and Inference}
\label{sec:design}
\subsection{Panels and statistical unit}
\begin{table}[tb]
\centering\small
\caption{Panel contracts. A block is one class--generation-seed pair. Conditional
replicates, anchors, coordinates and output images are not additional independent
samples. Overlap is assessed within a checkpoint, never by matching Tiny and
Small numeric block identifiers.}
\label{tab:panels}
\begin{tabular}{@{}p{0.37\textwidth}cccc@{}}\toprule
Panel & CFG & Blocks & Anchors & Replicates\\\midrule
P2P Tiny W8 primary & 3 & 48 & 5 & 8\\
P3S Small W8 primary & 3 & 40 & 2 & 8\\
P3S Small W8 sensitivity & 1 & 16 & 2 & 4\\
P2P Tiny W4 stress & 3 & 24 & 2 & 4\\
P2P matched timing & 3 & 24 & 2 & 2\\
P3S matched timing & 3 & 40 & 2 & 2\\\bottomrule
\end{tabular}
\end{table}
Within each panel, target $0.3$ and $0.5$ use the same block identifiers.
CFG1 uses a 16-block subset of the Small primary panel; W4 uses a 24-block
subset of the Tiny W8 panel. They are sensitivity/stress comparisons, not
independent extra replications. The compact files permit checking panel IDs,
counts and matching. They omit a complete raw class-ID/base-seed manifest, so
this reanalysis does not independently certify class--seed uniqueness or
predecessor disjointness. P2P compact CSVs omit CFG metadata; CFG3 is sourced
from the preserved protocol rather than independently filtered from those rows. Repeated seeds within a class are not automatically
pseudoreplication; the block bootstrap assumes blocks are independent for the
chosen class--seed population, with no separately assessed class-cluster effect.
A class-cluster reanalysis is unavailable without the true class manifest.

\subsection{Selection history and matched timing}
P2P follows the P2N observations, after which the $0.5$ versus $0.3$ contrast was
specified and the execution records describe new blocks. P3S records 40 new
blocks and a conservative power plan using half the P2P paired effect, candidates
32/40/48, and a 90\% target; the archived audit reports 40 as the first adequate
candidate. These are preserved historical planning records. Public YAML and
embedded timestamps alone do not provide independent external proof that the
analysis was registered before outcome access.

Matched timing retains the same seed coordinate and token substitution across
two anchor contexts according to the original execution logs. Its purpose is
control of seed identity. State, remaining horizon and coordinate context still
vary, so it is not a pure causal time intervention. Public compact matched files
lack the full token/coordinate identity needed to revalidate that construction
independently. Its score contrast is reproducible, with $n=24$ for Tiny and
$n=40$ for Small, and is kept separate from natural-shock results.

\subsection{Legacy inference and post-hoc supplement}
\label{sec:stats}
The headline conditional-LPIPS and incidence-weighted-LPIPS intervals retain the
original 20,000-resample paired block bootstrap, semantic SHA-256 label seeds,
and NumPy \texttt{higher} quantiles. Existing gate intervals for DINOv2 and
matched LPIPS use that same original rule. Each original contrast resamples
paired block differences, preserving the two anchors. Its semantic seed is
contrast-specific; the new supplement instead shares block indices across
metrics within each panel. Legacy reaggregation
checks the eight stored contrasts and two stored OOF MAE ratios; checking both
CI endpoints is stronger than the previous public script's lower-only checks.

A separately labeled \emph{post-hoc robustness supplement} uses 100,000 paired
block resamples and \texttt{linear} quantiles. Stable SHA-256 label-to-seed
mapping, RNG, resample count, input hashes, Python/NumPy/SciPy versions and
exact outputs are stored in the generated audit. Seeds are not based on
process-dependent Python \texttt{hash()}. Different seeds can shift Monte Carlo
endpoints; no endpoint is forced to match a comparison report.
Ratios and $\Msplit$ are recalculated inside each bootstrap sample.
Paired t intervals, leave-one-block-out influence, trimmed means and blockwise
signs are diagnostics, not randomized-treatment causal tests.

The supplement also reports 99.1667\% marginal intervals for the six
$H,\Meq,R$ contrasts in the \emph{two primary CFG3 W8 panels only}.
Bonferroni accounting for this six-endpoint family does not correct all
post-hoc hypothesis selection in the research program. CFG1, W4, matched timing
and interactions are outside that family. Conditional rollout counts are only
eight (primary), four (sensitivity/stress) or two (matched) per block/anchor;
block intervals do not separately identify all within-block Monte Carlo error.

\section{Results}
\label{sec:results}
\subsection{Primary contrast: lower batch incidence and higher discrepancy}
\label{sec:primary}
Table~\ref{tab:main_panels} gives the two primary panels. Both show
$H\downarrow$, $\Meq\uparrow$, $R\uparrow$ from target $0.3$ to $0.5$.
Legacy equal-block conditional LPIPS intervals are approximately
$[0.01719,0.02287]$ for Tiny and $[0.01714,0.02381]$ for Small;
incidence-weighted intervals are approximately $[0.001368,0.002066]$ and
$[0.001499,0.002477]$. These are within-checkpoint paired contrasts.
The exact source-linked values and sample counts are generated below.
\begin{table}[tb]
\centering\small
\setlength{\tabcolsep}{3pt}
\caption{Primary CFG3 W8 block means and paired contrasts, $\Delta=.5-.3$.
$\Meq$ is equal-block conditional LPIPS and $R=\operatorname{mean}_b(hm)$.
Conditional-LPIPS and $R$ headline intervals use the original 20k/higher rule;
$H$ inference is supplied only in the post-hoc supplement.}
\label{tab:main_panels}
\input{\revisiongenerated/main_panels.tex}
\end{table}
\begin{figure}[tb]
\centering\includegraphics[width=\textwidth]{\revisiongenerated/main_panels.pdf}
\caption{Separate axes for $H$, $\Meq$ and $R$ in the primary CFG3 W8 panels;
block counts are panel-specific. Figure intervals use the labeled post-hoc
100k/linear paired-block rule. Directions concern this one-shock intervention,
not persistent approximation or image quality.}
\label{fig:main}
\end{figure}
The three directions hold simultaneously in 43/48 Tiny and 39/40 Small blocks.
This is a descriptive blockwise pattern, not a theorem. DINOv2 supports a
positive equal-block conditional contrast in the original hypotheses, with
absolute effects differing by checkpoint. Tiny's full five-anchor $R$ profile
is not strictly monotone: the $0.1\rightarrow0.3$ contrast is about
$-0.000054$ with a post-hoc interval spanning zero. The final-flush anchor must
therefore remain visible in the full-anchor supplement.

\subsection{Aggregation and attribution}
\label{sec:attribution}
The equal-block and split-population conditional means differ because $h$ and
$m$ covary. Tiny $\Msplit$ is about $0.010378\rightarrow0.030267$;
Small $\Msplit$ is about $0.010716\rightarrow0.032086$.
In Small, $H\Meq$ is about 3.13\% below $R$ at target $0.3$ and 4.14\%
below it at $0.5$. These are covariance effects, not a rounding discrepancy.
\begin{table}[tb]
\centering\small
\setlength{\tabcolsep}{3pt}
\caption{Aggregation check from block data. Population covariance uses $1/n$.
$\Msplit=R/H$ is not the historical equal-block conditional column.}
\label{tab:aggregation}
\input{\revisiongenerated/aggregation.tex}
\end{table}
For the $0.5-0.3$ contrast the exact symmetric arithmetic decomposition is
\begin{align}
 A_h&=\operatorname{mean}_b\left[(h_{b,.5}-h_{b,.3})
                    \frac{m_{b,.5}+m_{b,.3}}2\right],\\
 A_m&=\operatorname{mean}_b\left[(m_{b,.5}-m_{b,.3})
                    \frac{h_{b,.5}+h_{b,.3}}2\right],
 &\hspace{1em}\Delta R=A_h+A_m.
\end{align}
Apply the identity per block and then average. Replacing the block expressions
by products of panel averages omits covariance terms. Tiny contributions are
approximately $-0.001743$ and $+0.003457$; Small contributions are approximately
$-0.001511$ and $+0.003477$. This accounting is not a causal mediation analysis.
\begin{table}[tb]
\centering\small
\setlength{\tabcolsep}{3pt}
\caption{Blockwise symmetric attribution and simultaneous blockwise directions.
The identity is checked in the generated audit; the components are descriptive,
not independent causal effects.}
\label{tab:attribution}
\input{\revisiongenerated/attribution.tex}
\end{table}

\subsection{Schedule and local-TV diagnostics}
\label{sec:schedule}
The public expected seed count is $\sum_{u\in C}\TV(p_u,q_u)$, so mean
coordinate TV equals expected seed count divided by commit batch size $c$.
This average does not reconstruct the full coordinate TV distribution.
For $0.3\rightarrow0.5$, $c$ decreases from 14 to 8. Tiny's mean coordinate
TV is about $0.017015\rightarrow0.017232$ and its paired difference interval
includes zero; this is not an equivalence result. Small's mean coordinate TV
increases, approximately $0.015846\rightarrow0.016988$, with a positive
paired difference interval. Batch $h$ nevertheless falls in both panels.
\begin{table}[tb]
\centering\small
\setlength{\tabcolsep}{3pt}
\caption{Commit opportunities and mean coordinate TV, calculated from the public
expected-seed-count field. Intervals are post-hoc 100k/linear paired-block
intervals for $\Delta=.5-.3$.}
\label{tab:schedule_tv}
\input{\revisiongenerated/schedule_tv.tex}
\end{table}
\begin{figure}[tb]
\centering\includegraphics[width=\textwidth]{\revisiongenerated/schedule_tv.pdf}
\caption{Batch schedule and mean-coordinate TV displayed separately. Falling
batch split probability does not establish falling local quantization error.
Remaining horizon, state, positions and seed composition also change.}
\label{fig:schedule}
\end{figure}
The product split law depends on how many coordinates can split as well as
their TVs. Schedule/branching opportunity is therefore an explanation candidate,
not an identified sole cause. The current observations do not isolate pure
batch-size or time effects.

\subsection{Matched timing, sensitivity and saturated stress}
\label{sec:auxiliary}
Matched timing gives a separate equal-block LPIPS contrast of about $0.025088$
for Tiny ($n=24$) and $0.022235$ for Small ($n=40$).
\begin{table}[tb]
\centering\small
\setlength{\tabcolsep}{3pt}
\caption{Matched-timing score contrasts are separate from natural-shock $R$.
Tiny has 24 paired blocks, Small 40; each has two conditional replicates.}
\label{tab:matched}
\input{\revisiongenerated/matched.tex}
\end{table}
The official CFG1 Small sensitivity panel has 16 overlapping blocks. Its
$\Meq$ and $R$ contrasts have the same signs as CFG3; the paired CFG1-minus-CFG3
interaction is computed only on their common subset. Directional consistency
does not establish equal effect sizes. The Tiny W4 stress subset uses 24
blocks overlapping W8. Its $R$ is about $0.029407\rightarrow0.051431$,
while $H$ is near saturation ($0.9878\rightarrow0.9367$). It is an exploratory
stress comparison; paired W4-versus-W8 comparisons use only their common subset.
Neither subset is an independent non-saturated W8 replication.
\begin{table}[tb]
\centering\small\setlength{\tabcolsep}{3pt}
\caption{Exploratory paired differences in $\Delta=.5-.3$, restricted to common
within-checkpoint blocks. Same direction does not imply equal effect size.}
\label{tab:interactions}
\input{\revisiongenerated/interactions.tex}
\end{table}


\subsection{Prediction evidence and its reproducibility limit}
\label{sec:prediction}
The archived analyses fit leave-one-block-out ridge models. A seed-only model
uses anchor category, count, severity, surprisal and seed topology; adding
post-suffix descendant features yields stored OOF MAE ratios of approximately
0.5166 (Tiny) and 0.5440 (Small). The CPU reproduction can reaggregate stored
OOF errors. That is not an independent refit of the feature extraction and LOBO
training pipeline. Descendant information is measured after propagation;
it cannot be used as an online pre-terminal precision-selection feature.
No current result establishes improvement over a strong prospective policy
baseline at matched parameter bytes or measured cost.

\section{Negative Findings and Limitations}
\label{sec:limitations}
P05R's finite-profile influence envelope was conservative but empirically
vacuous in earlier regimes; P1M found little within-anchor incremental value.
P2A remained incomplete at its historical stopping point; P2B subsequently
closed the particular RNG/MASK/clamp audit questions in its own recorded stage.
The revision preserves these outcomes and does not retroactively relabel P2A
complete. Associations involving source age were unstable under adjustment,
and no source-age causal claim is made. No discontinuous phase transition was
fitted or tested.

The two checkpoints share an architecture family, VQ decoder and scheduler.
The public data do not establish architecture, scheduler or coupling universality.
A split hazard for persistent paths is different from these anchor-conditioned
batch probabilities. LPIPS and DINOv2 measure specified output discrepancies;
no human similarity, preference or quality study is included. Weight-only fake
quantization supplies no packed-kernel speed or memory measurements.

The block bootstrap assumes independent class--seed blocks and addresses the
empirical block population. Compact IDs do not reveal all class/seed provenance,
predecessor overlap or potential class-level dependence. Repeating multiple seeds
of a class is not by itself proof of pseudoreplication, and the original primary
analysis already used block resampling. A class-cluster check needs the true
manifest. Limited conditional rollout counts create additional Monte Carlo
uncertainty. Historical audits are evidence about the runs they inspected;
this CPU revision does not independently rerun their full model/metric pipeline,
rehash excluded weights or validate every original token identity. Stored $h$
and LPIPS/DINO scores are reaggregated, not independently reconstructed from
per-coordinate categorical marginals or terminal images absent from these tables.

The historical per-tensor plan differs from frozen per-channel implementation
source (Section~\ref{sec:quantizer}); that deviation is now explicit. Exact
historical runtime scale/weight tensor replay has not been performed. The theoretical influence
envelope requires global coverage unavailable from finite neural local profiles,
and no valid token-bound-to-LPIPS transfer constant is supplied.

\section{Unexecuted Research Plan and Falsification Criteria}
\label{sec:future}
Every experiment in this section is a proposal. No new result or compute budget
is implied by its inclusion.

\paragraph{P1: schedule and split opportunities.}
First test the same model in a $2\times2$ design: the existing batch sequence
versus $c=8$ for each of 32 steps, and Halton order versus a fixed permutation.
Changing the schedule also changes prefix states, so it is not a pure causal
intervention on $c$. A complementary diagnostic uses one identical anchor
state and nested sets of $c=1,4,8$ coordinates: compare $p$ with $q$ on those
positions, use common $p$ on all others, and then common $q$ in the suffix.
This is a diagnostic intervention distinct from the full natural quantization
transition. Freeze comparisons before evaluation. If equal-$c$ removes the
incidence reversal, narrow the interpretation to schedule-dependent opportunities;
if conditional/weighted differences vanish after seed count and severity are
controlled, do not claim residual timing sensitivity.

\paragraph{P2: a different discrete generation structure.}
A first candidate is single-token AR generation such as LlamaGen, subject to
separate verification of official checkpoint, tokenizer and runtime feasibility;
no GPU memory or time is guaranteed. Define the same shock/common-suffix contract
with $c=1$. KV cache is part of state. A pure token-mediated arm must reconstruct
a common-model cache from the canonical prefix and be distinguished from a
persistent-cache arm. Report within-model paired effects rather than pooling
raw LPIPS across models. Failure to repeat $H\downarrow$ would refute universal
sign-pattern replication, even if incidence/consequence separation remains
useful. Continuous diffusion or MAR with threshold-defined onset introduces a
different estimand.

\paragraph{P3: non-saturated precision and coupling robustness.}
After establishing W8 granularity, select W6 or one alternative calibration
method on development data and freeze it before held-out evaluation. W4
saturation is not primary generalization evidence. For another coupling,
validate marginals and self-cases and report how $h$ and paired discrepancy
change. If effect signs or policy conclusions fail across predeclared valid
couplings, report coupling dependence rather than universality.

\paragraph{Separate practical-value test.}
Retain any existing layer-by-time plan, but distinguish persistent W8 in one
module, a one-time shock, and restoring one module to FP32 from an all-W8
baseline. Do not assume quantizing one module and restoring one from all-W8 are
symmetric or additive. Compare local MSE, logit KL/TV, $h$ alone, seed count and
severity, simple early protection and same-cost random selection. Freeze a
policy on development blocks, then evaluate actual restoration gain on entirely
new class/seed blocks at matched parameter bytes or measured cost. Exclude
post-hoc descendant features from online predictor inputs. If the method offers
no added gain over these baselines, narrow the practical contribution.
Separate cost/variance pilots from confirmation and choose sample size from
effect size, power and multiplicity, without guaranteeing that a particular
block count suffices. Any human study must distinguish blinded paired similarity
from quality/preference and use representative samples; none is present here.

\section{Conclusion}
FDCA specifies how to join a natural approximation-induced token split to its
conditional output discrepancy. Under a fixed Halton schedule, coordinatewise
coupling and a one-step FP32-versus-W8/common-W8-suffix intervention, two CFG3
checkpoint panels support lower batch incidence but higher equal-block
conditional and incidence-weighted LPIPS at target $0.5$ than at $0.3$.
The weighted result is an average of blockwise products, and batch incidence
also depends on commit opportunities. The theory supplies finite monotone-decoder
bounds and an expectation identity, not a universal trend or a perceptual quality
guarantee. Prospective prediction, precision selection and new-architecture
validation remain open experiments.

\section*{Data, Code, Version and Citation}
The public repository is
\url{https://github.com/Munsik-Kim/fdca-visual-token}.
Archived v1.0.0 remains associated with
\href{https://doi.org/10.5281/zenodo.21965747}{10.5281/zenodo.21965747}.
This GitHub revision candidate is not a new release or Zenodo deposit. The author,
ORCID and license scopes are preserved: original code is MIT; paper, figures,
tables and included audit documentation are CC BY 4.0. Third-party weights,
data and upstream code retain their own rights and are not redistributed.

The reanalysis script generates the numerical JSON/CSV, manuscript table inputs
and vector/raster figures from the original compact tables. Its audit records
input SHA-256 hashes, environment, RNG mapping and resampling rules. The
version-aware paper build compiles this source with its updated bibliography;
legacy frozen-paper verification remains distinct. Revision notes identify
corrections, unresolved items and historical evidence boundaries.

\appendix
\section{Post-hoc Robustness Supplement}
\label{app:robustness}
These analyses were added after the original experiments. They are not the
original preregistered analysis. All paired differences are $0.5-0.3$ within
one panel. Table~\ref{tab:posthoc} uses 100,000 paired block bootstrap draws
and linear quantiles, with ratio and split-population conditional means
recalculated per resample. It is intentionally separate from the original
20k/higher headline intervals.
\begin{table}[p]
\centering\small
\setlength{\tabcolsep}{3pt}
\caption{Post-hoc natural-shock contrasts, with panel-specific $n$. CFG1 and W4
are overlapping sensitivity/stress subsets. The precise fixed RNG and input
hashes are in the generated audit.}
\label{tab:posthoc}
\input{\revisiongenerated/posthoc_contrasts.tex}
\end{table}
\begin{table}[p]
\centering\small
\setlength{\tabcolsep}{3pt}
\caption{Robustness diagnostics and the explicitly limited six-contrast
Bonferroni family where indicated. Paired t, leave-one-block-out and trimmed
summaries are descriptive sensitivity checks, not randomized causal inference.}
\label{tab:robustness}
\input{\revisiongenerated/robustness.tex}
\end{table}
\begin{figure}[tbp]
\centering\includegraphics[width=\textwidth]{\revisiongenerated/paired_contrasts.pdf}
\caption{Post-hoc paired contrasts, $0.5-0.3$, resampling complete block bundles.
The displayed interval rule is 100k/linear; axes and panel counts distinguish
probability from LPIPS metrics.}
\label{fig:contrasts}
\end{figure}
\begin{figure}[tbp]
\centering\includegraphics[width=\textwidth]{\revisiongenerated/tiny_all_rho.pdf}
\caption{Full Tiny CFG3 W8 anchor profile ($n=48$). All five target fractions
are shown, including the final-flush $0.1$ anchor. The mean of blockwise
products $R$ does not increase at every adjacent anchor. This full-anchor descriptive figure shows means without intervals; the
selected $0.3/0.5$ region is shaded, and adjacent contrasts are supplied in the
generated tables.}
\label{fig:allrho}
\end{figure}

\FloatBarrier
\section{Historical Evidence and Proof Ledger}
The original program progressed through T0 exact finite-state checks, P0 sampler
inspection, P05R influence measurement, P1L controlled shocks, P1M incremental
analysis, P2N natural seeds, P2A's incomplete audit, P2B's targeted closure,
P2P perceptual measurements and P3S second-checkpoint replication. This order
matters for selection: the P2P contrast followed earlier P2N findings; later
new-block records are not equivalent to an externally timestamped preregistration.

The archived P3S report distinguishes 52 executed pytest categories from 60
generated audit assertions. It reports 1,429 run-manifest files, whereas the
Full package contains 1,431 entries including two bookkeeping files. Its 1,040
paired-grid files include primary, matched, controls and the CFG1 subset; this
count is not a count of independent blocks. The audit did not independently
rehash excluded checkpoint bytes or rerun a Parquet pipeline in its environment.
Likewise, current OOF MAE reaggregation is not independent model training.

The proof ledger in \nolinkurl{docs/PROOF_LEDGER_v1.0.1.md} traces each theorem's
assumptions and proof steps. The stable-set theorem is verified for its stated
nonnegative bounded additive costs. The full-path decomposition is verified as
a conditional-expectation identity with survival and zero no-split cost. The
DAG result is verified conditional on the stated global hybrid closure and
coordinate coupling; empirical local neural certificates and token-to-LPIPS
transfer remain unproved. The ledger does not classify a familiar-looking
statement as verified without those proof connections.

\bibliography{references}
\end{document}
